\subsection{Arbitrary independent blocks, including hard constraints}
\label{sec:threshold-positive-blocks}

\begin{theorem}[Block-product full AA without nondegeneracy]
\label{thm:threshold-block-products}
Let the input coordinates be partitioned into blocks of size at most
\(k\), and let \(\mu\) be an arbitrary product of laws on these
blocks. The block laws may depend on dimension and may assign
probability zero to some atoms. For \(T\ge1\) and \(0<e<1\), set
\[
 \eta=\frac{e^2}{64T^2},\qquad
 b=k+\left\lceil\log_2(1/\eta)\right\rceil.
\]
Every \(T\)-query acceptance function admits a deterministic simulator
of mean error at most \(e\), using at most
\[
 S_B(T,e;k):=k\,s_U(2bT,e/2)
\]
queries. In particular, a uniform bound on block size implies full AA,
uniformly over all choices of block laws and temperatures.
\end{theorem}

\begin{proof}
For a block law \(p\) on at most \(2^k\) atoms, put \(N=2^b\).
Give each atom \(u\) the integer count
\(c_u=\lceil(1-\eta)Np(u)\rceil\). Since
\(\sum_uc_u\le(1-\eta)N+2^k\le N\), distribute any remaining
counts arbitrarily. The resulting dyadic law \(p'\) satisfies
\(p'(u)\ge(1-\eta)p(u)\), so it has the form
\[
 p'=(1-\eta)p+\eta r
\]
for another probability law \(r\). A deterministic map of \(b\)
fair bits generates \(p'\) exactly.

Perform this construction independently in each block. Couple
\(X\sim\mu\) to the product law \(Y\) by keeping each entire block
with probability \(1-\eta\), otherwise replacing it by an independent
sample from its \(r\). Thus, conditional on \emph{all} of \(X=x\),
each coordinate changes with probability at most \(\eta\).

Compose the block gadgets with \(Q\). Each virtual input query costs
at most \(2b\) fair-source queries by reading the source block,
computing the requested output bit, and uncomputing. Uniform AA gives
a source tree with error \(e/2\). To supply its source bits on input
\(X\), query the entire physical block when that source block is first
visited, generate its coupled \(Y\)-block, and sample the uniform
source block conditional on that \(Y\)-block. Cache the result.
Different blocks are independent, so this is the correct conditional
source oracle, with at most \(k\) physical queries per first visit.
The transfer proposition gives error
\(e/2+4T\sqrt\eta=e\), followed by derandomization.
\end{proof}

The case \(k=1\) covers products with arbitrarily varying coordinate
biases, including deterministic bits, with no minimum bias assumption.
The case \(k=3\) includes products of laws supported on
\(\{+++ ,---\}\), for which the flip-influence property can fail even
at degree three. Thus the absence of a nondegeneracy assumption is
specific to the full query-simulation target. For a Hamiltonian that
is a sum of disjoint bounded-size components, no finite-temperature
full-AA transition is possible under this theorem; even hard-constraint
limits of the block laws are covered.
